% Lösungen Lineare Funktionen (zusammenbau v0.4, Heft)
\einheitenkopf{Lineare Funktionen}

\erg{1}{a) $-5$, $-2$, $1$, $4$ \quad b) Start $(0|4)$, 1 nach rechts, 3 nach unten: Gerade durch $(0|4)$ und $(1|1)$}
\erg{2}{a) Start $(0|3)$, 1 nach rechts, 4 nach unten: Gerade durch $(0|3)$ und $(1|-1)$ \quad b) Start $(0|-4)$, 1 nach rechts, 3 nach unten: Gerade durch $(0|-4)$ und $(1|-7)$ \quad c) Start $(0|-3)$, 1 nach rechts, 5 nach oben: Gerade durch $(0|-3)$ und $(1|2)$ \quad d) Start $(0|-4)$, 1 nach rechts, 6 nach oben: Gerade durch $(0|-4)$ und $(1|2)$ \quad e) Start $(0|5)$, 1 nach rechts, 3 nach unten: Gerade durch $(0|5)$ und $(1|2)$ \quad f) Start $(0|-2)$, 1 nach rechts, 4 nach unten: Gerade durch $(0|-2)$ und $(1|-6)$ \quad g) Start $(0|5)$, 1 nach rechts, 2 nach oben: Gerade durch $(0|5)$ und $(1|7)$ \quad h) Start $(0|-3)$, 1 nach rechts, 4 nach oben: Gerade durch $(0|-3)$ und $(1|1)$}
\erg{3}{a) $n = 3$ \quad b) $m = -2$ (1 nach rechts, 2 nach unten)}
\erg{4}{a) f: $y = -3x + 3$, g: $y = 2x + 3$, h: $y = -3$ \quad b) f: $y = -4x + 4$, g: $y = 2x + 4$, h: $y = -1$ \quad c) Anstieg: f; parallel zur x-Achse: h \quad d) Anstieg: g; parallel zur x-Achse: h \quad e) g (von links nach rechts gelesen geht sie nach unten) \quad f) f (von links nach rechts gelesen geht sie nach unten) \quad g) $(0|-3)$ \quad h) $(0|3)$ \quad i) A (eine Gerade ohne Knick) \quad j) C (eine Gerade ohne Knick)}
\erg{5}{a) $f(4) = 2 \cdot 4 + 5 = 13$ \quad b) $f(-3) = 4 \cdot (-3) - 1 = -13$}
\erg{6}{a) Gerade durch $(0|6)$ und $(1|2)$; $0 = -4x + 6$, Nullstelle $x = 1{,}5$ \quad b) Gerade durch $(0|3)$ und $(1|-3)$; $0 = -6x + 3$, Nullstelle $x = 0{,}5$ \quad c) $0 = -0{,}4x + 60$, $x = 60 : 0{,}4 = 150$; nach 150 Minuten \quad d) $0 = -0{,}8x + 100$, $x = 100 : 0{,}8 = 125$; nach 125 Minuten \quad e) $x = -2$ (dort schneidet g die x-Achse) \quad f) $x = 4$ (dort schneidet g die x-Achse) \quad g) Schnittpunkt $(2|3)$; ja, zwei Punkte: der zweite liegt außerhalb des Bildes bei $x = -3$ \quad h) Schnittpunkt $(1|6)$; ja, zwei Punkte: der zweite liegt außerhalb des Bildes bei $x = -5$ \quad i) $y = \frac{1}{3} \cdot (-9) - 2 = -5$ \quad j) $y = \frac{3}{4} \cdot (-8) + 2 = -4$ \quad k) Ja: $f(-4) = 5$ \quad l) Nein: $f(-2) = 2$, nicht $4$ \quad m) Die Gerade f schneidet die y-Achse in $P(0|6)$. Die Gerade f ist parallel zur Geraden $y = 2x$. \quad n) richtig sind „Die Gerade f fällt“ und die Aussage zum Punkt R \quad o) $(-3|16)$, denn $f(-3) = 20$ \quad p) $(-2|-6)$, denn $f(-2) = 14$}
\erg{7}{a) $m = 2$, $n = -3$: $f(x) = 2x - 3$ \quad b) $m = (0 - 6) : (3 - 1) = -3$, $n = 9$: $f(x) = -3x + 9$}
\erg{8}{a) Gerade durch A und B; erste Aussage falsch, zweite wahr; $m = -1{,}5$, $n = 5$: $f(x) = -1{,}5x + 5$ \quad b) Gerade durch A und B; erste Aussage falsch, zweite wahr; $m = 0{,}5$, $n = -3$: $f(x) = 0{,}5x - 3$ \quad c) $m = (440 - 800) : (40 - 10) = -12$; $800 = -12 \cdot 10 + n$, $n = 920$: $h(t) = -12t + 920$ \quad d) $m = (500 - 260) : (10 - 4) = 40$; $260 = 40 \cdot 4 + n$, $n = 100$: $V(t) = 40t + 100$ \quad e) $m = (0 - (-4)) : (3 - (-3)) = \frac{2}{3}$; aus $D$: $0 = \frac{2}{3} \cdot 3 + n$, $n = -2$ – die Gleichung stimmt \quad f) $m = (2 - 5) : (4 - (-2)) = -\frac{1}{2}$; aus $F$: $2 = -\frac{1}{2} \cdot 4 + n$, $n = 4$ – die Gleichung stimmt}
\erg{9}{a) Anfangswert 45 € ($n = 45$), Änderung 38 € je Stunde ($m = 38$) \quad b) $y = 3x + 8$ \quad c) $m = 6$, $n = 10$: $y = 6x + 10$}
\erg{10}{a) Angebot 1: $4 \cdot 32 + 120 \cdot 0{,}30 = 164$ €; Angebot 2: $95 + 420 \cdot 0{,}12 = 145{,}40$ €; Angebot 2 ist günstiger; $y = 0{,}12x + 95$ \quad b) Angebot 1: $3 \cdot 29 + 130 \cdot 0{,}40 = 139$ €; Angebot 2: $110 + 380 \cdot 0{,}10 = 148$ €; Angebot 1 ist günstiger; $y = 0{,}10x + 110$ \quad c) 240 km: Nord $150 + 1{,}20 \cdot 240 = 438$ €, Süd $1{,}80 \cdot 240 = 432$ €, Süd günstiger; 320 km: Nord $150 + 1{,}20 \cdot 320 = 534$ €, Süd $1{,}80 \cdot 320 = 576$ €, jetzt Nord günstiger (gleich teuer bei 250 km) \quad d) 260 km: Adler $90 + 0{,}80 \cdot 260 = 298$ €, Biber $1{,}10 \cdot 260 = 286$ €, Biber günstiger; 320 km: Adler $90 + 0{,}80 \cdot 320 = 346$ €, Biber $1{,}10 \cdot 320 = 352$ €, jetzt Adler günstiger (gleich teuer bei 300 km) \quad e) erster Sachverhalt: $y = 0{,}25x + 1$ (Cent in Euro umrechnen), zweiter: $y = x + 25$ \quad f) erster Sachverhalt: $y = 0{,}60x + 2$ (Cent in Euro umrechnen), zweiter: $y = 2x + 60$ \quad g) $m = 45$, $n = 640$: $y = 45x + 640$; aus $45x + 640 = 1\,500$ folgt $x \approx 19{,}1$, also nach 20 Monaten \quad h) $m = 70$, $n = 1\,150$: $y = 70x + 1\,150$; aus $70x + 1\,150 = 2\,000$ folgt $x \approx 12{,}1$, also nach 13 Monaten \quad i) $m = 1{,}2$, $n = 180$: $y = 1{,}2x + 180$ \quad j) $m = 2{,}4$, $n = 250$: $y = 2{,}4x + 250$ \quad k) $830 + 12 \cdot 35 = 1\,250$ € \quad l) $2\,340 + 6 \cdot 85 = 2\,850$ Bäume \quad m) zu Beginn $300$ l (30 l mehr als nach 5 min); nach 40 min: $300 - 8 \cdot 30 = 60$ l \quad n) zu Beginn $20$ \% (6 Prozentpunkte weniger als nach 10 min); nach 90 min: $20 + 9 \cdot 6 = 74$ \% \quad o) Blitz: II, Flott: I; II beginnt bei 40 € auf der y-Achse – das ist der Grundpreis \quad p) Hoch: I, Weit: II; I beginnt bei 30 € auf der y-Achse – das ist die Aufnahmegebühr \quad q) erstes Angebot: A (bis zu den Freikilometern gleichbleibend, dann steigend); zweites Angebot: C (beginnt bei der Einmalzahlung, steigt gleichmäßig) \quad r) erstes Angebot: D (bis zu den Freikilometern gleichbleibend, dann steigend); zweites Angebot: B (beginnt bei der Einmalzahlung, steigt gleichmäßig)}
